\chapter{Requirements Specification and Analysis}
\chaptertoc

\section{Introduction}
This chapter formalizes system requirements and the main design choices that structure the project implementation. It defines actors, functional and non-functional requirements, product backlog orientation, sprint planning, data model foundations, project environment, and software architecture.

\section{Requirements Specification}

\subsection{Identification of Actors}
The system involves three primary actors:
\begin{itemize}
    \item \textbf{Administrator}: configures sources, thresholds, and integration parameters.
    \item \textbf{Operator}: monitors live queues and reacts to alerts during operations.
    \item \textbf{Manager}: analyzes historical trends and uses indicators for planning decisions.
\end{itemize}

\subsection{Functional Requirements}
The system must provide the following core functionalities:
\begin{itemize}
    \item real-time person detection and tracking,
    \item configurable queue zone definition,
    \item estimation of arrival rate ($\lambda$), service rate ($\mu$), and wait time ($W$),
    \item uncertainty quantification for queue metrics,
    \item real-time web dashboard updates and alerts,
    \item RTSP and ONVIF camera integration,
    \item automated webhook and Telegram notifications,
    \item historical data logging and analytics.
\end{itemize}

\subsection{Non-functional Requirements}
The main quality constraints are:
\begin{itemize}
    \item \textbf{Performance}: maintain real-time behavior and low alert latency.
    \item \textbf{Scalability}: support multiple feeds and independent feed lifecycle control.
    \item \textbf{Reliability}: tolerate transient source failures and integration interruptions.
    \item \textbf{Usability}: provide clear web workflows for monitoring and configuration.
    \item \textbf{Security}: protect integration channels and service endpoints.
    \item \textbf{Maintainability}: keep modular code organization and testable interfaces.
\end{itemize}

\section{Product Backlog}
The product backlog was prioritized by business value and technical dependency. For implementation governance, items were grouped into three domains:
\begin{itemize}
    \item \textbf{Backend domain}: perception pipeline, queue analytics, API and stream interfaces.
    \item \textbf{Web domain}: dashboard UX, live monitoring integration, zone editing workflows.
    \item \textbf{Automation domain}: n8n workflow integration, alert routing, webhook reliability.
\end{itemize}
Early sprints concentrated on backend foundations to reduce downstream integration risk.

\section{Global Use Case Diagram}
\begin{figure}[ht]
    \centering
    \fbox{\parbox[c][6cm][c]{0.85\textwidth}{\centering Placeholder for global use case diagram}}
    \caption{Global use case diagram of the queue wait-time estimation system.}
    \label{fig:global_use_case}
\end{figure}

\section{Sprint Scope Overview}
The project implementation is organized into eight sprints. Table~\ref{tab:sprint_plan} summarizes sprint scope and expected outcomes.

\begin{table}[ht]
    \centering
    \caption{Project sprint scope overview}
    \label{tab:sprint_plan}
    \begin{tabularx}{\textwidth}{l X}
        \toprule
        Sprint & Main outcome \\
        \midrule
        Sprint 1 & Detection, tracking, and queue-zone perception foundation \\
        Sprint 2 & Queue analytics and uncertainty quantification \\
        Sprint 3 & Camera connectivity and backend API exposure \\
        Sprint 4 & Web design system and frontend scaffolding \\
        Sprint 5 & Frontend integration and live monitoring \\
        Sprint 6 & n8n automation and Telegram alerting \\
        Sprint 7 & Evaluation and optimization \\   
        Sprint 8 & Final consolidation, reporting, and delivery \\
        \bottomrule
    \end{tabularx}
\end{table}

\section{Data Modeling}
The data model combines runtime monitoring entities and persistence entities. The main objects are:
\begin{itemize}
    \item \textbf{VideoFeed}: source identity, connectivity, runtime status, and configuration.
    \item \textbf{QueueMetrics}: timestamped indicators (people in zone, $\lambda$, $\mu$, $W$, stability).
    \item \textbf{Alert}: severity, threshold context, and notification payload.
    \item \textbf{Establishment} and \textbf{Caisse}: business metadata used to organize feeds.
\end{itemize}
These entities support both live supervision and historical analysis pipelines.

\section{Project Environment}

\subsection{Hardware and Software Environment}
The development workflow uses a Python backend environment, a TypeScript frontend environment, and containerized automation services. The project is designed to run on standard development workstations and can exploit GPU acceleration when CUDA-compatible hardware is available.

\subsection{Technologies Used}
The selected technology stack is:
\begin{itemize}
    \item \textbf{Vision pipeline}: YOLO26, ByteTrack, OpenCV \cite{ultralytics_docs,bytetrack2022,opencv_docs}.
    \item \textbf{Queue analytics}: NumPy, Pandas, SciPy, M/M/1 modeling \cite{gross_queueing}.
    \item \textbf{Backend API}: FastAPI and Pydantic v2.
    \item \textbf{Frontend}: React, TypeScript, Vite, Tailwind CSS.
    \item \textbf{Automation}: n8n workflows with Telegram notifications.
    \item \textbf{Persistence}: CSV metrics logging and SQLite metadata.
\end{itemize}

\section{Software Architecture}

\subsection{Overall Architecture}
The system architecture is organized in four layers:
\begin{enumerate}
    \item \textbf{Video acquisition}: local files, RTSP streams, and ONVIF discovery.
    \item \textbf{Perception}: detection, tracking, and queue-zone filtering.
    \item \textbf{Analysis}: queue metrics estimation and uncertainty modeling.
    \item \textbf{Output and integration}: web dashboard, webhook dispatch, and alert automation.
\end{enumerate}

\subsection{Backend Architecture}
The backend exposes REST and WebSocket interfaces for feed lifecycle control and live metrics. It includes modules for detector orchestration, tracking, queue analysis, threshold detection, webhook delivery, and camera-source management (RTSP and ONVIF). The API contract supports onboarding, runtime control, snapshots, streams, metadata, and integration diagnostics.

\subsection{Frontend Architecture}
The frontend follows a component-based architecture with typed API client mapping, feed-oriented dashboard views, and WebSocket-driven state updates. It integrates live stream rendering, per-feed controls, and alert visualization while preserving separation between presentation components and data-access logic.

\section{Conclusion}
This chapter established the requirements baseline and architecture model used throughout implementation. The next chapters present sprint-by-sprint development aligned to this specification and analysis framework.
